%
% Ableitung einer Potenzreihe
%
\subsection{Potenzreihen
\label{buch:holomorph:holomorph:subsection:potenzreihen}}
In Abschnitt~\ref{buch:komplex:funktionen:subsection:potenzreihen}
wurden einige aus der reellen Analysis wohlbekannte Funktionen
durch ihre Potenzreihe definiert.
In diesem Abschnitt soll gezeigt werden, dass alle diese Funktionen
holomorph sind.

%
% Ableitung einer Potenzreihe
%
\subsubsection{Ableitung einer Potenzreihe}
In der reellen Analysis wird gezeigt, dass eine Funktionenreihe gliedweise
abgeleitet werden darf, wenn die abgeleitete Reihe weiterhin konvergiert.

\begin{satz}[Ableitung einer Potenzreihe]
\label{buch:holomorph:holomorph:satz:potenzreihe-ableitung}
Sei eine Funktion durch die Potenzreihe
\index{Potenzreihe}%
\[
f(z)
=
\sum_{k=0}^\infty a_k z^k
\]
gegeben.
Dann ist
\[
f'(z)
=
\sum_{k=1}^\infty ka_k z^{k-1}
\]
eine konvergente Potenzreihe mit dem gleichen Konvergenzradius.
\index{Konvergenzradius}%
\end{satz}

\begin{proof}
Es reicht, den Konvergenzradius mit der Cauchy-Hadamard-Formel zu
bestimmen.
Dabei kann verwendet werden, dass der Konvergenzradius $\varrho$
der Potenzreihe für $f(z)$ bekannt ist und den Wert
\[
\frac{1}{\varrho}
=
\limsup_{n\to\infty} \!\sqrt[n]{|a_n|}
\]
hat.
Für die Koeffizienten der abgeleiteten Reihe folgt
\begin{align}
\limsup_{n\to\infty} \!\sqrt[n]{|na_n|}
&=
\limsup_{n\to\infty}
\!\sqrt[n]{n}
\!\sqrt[n]{|a_n|}.
\label{buch:holomorph:holomorph:satz:ableitungkonvergenzradius}
\end{align}
Wegen
\[
\frac{\log n}{n} \to 0\qquad\text{für $n\to\infty$}
\]
konvergiert die Folge $\!\sqrt[n]{n}$ gegen
\[
\lim_{n\to\infty}
\!\sqrt[n]{n}
=
\lim_{n\to\infty}
\exp \biggl(
\frac{1}{n}\log n
\biggr)
=
\exp \biggl(
\lim_{n\to\infty}
\frac{1}{n}\log n
\biggr)
=
e^0
=
1.
\]
Der Grenzwert 
\eqref{buch:holomorph:holomorph:satz:ableitungkonvergenzradius}
wird damit zu
\[
\limsup_{n\to\infty}
\!\sqrt[n]{n}
\!\sqrt[n]{|a_n|}
=
\lim_{n\to\infty}
\!\sqrt[n]{n}
\cdot
\limsup_{n\to\infty}
\!\sqrt[n]{|a_n|}
=
\limsup_{n\to\infty}
\!\sqrt[n]{|a_n|}
=
\frac{1}{\varrho}.
\]
Der Konvergenzradius der abgeleiteten Reihe ist der gleiche.
\end{proof}

\begin{satz}[Stammfunktion einer Potenzreihe]
\label{buch:holomorph:holomorph:satz:potenzreihe-stammfunktion}
\index{Stammfunktion}%
Sei eine Funktion durch die Potenzreihe
\[
f(z)
=
\sum_{k=0}^\infty a_k z^k
\]
gegeben.
Dann ist
\begin{equation*}
F(z)
=
\sum_{k=0}^\infty \frac{1}{k+1}a_k z^{k+1}
\end{equation*}
eine konvergente Potenzreihe mit dem gleichen Konvergenzradius.
$F(z)$ ist zudem eine Stammfunktion von $f$: $F'(z)=f(z)$.
\end{satz}

\begin{proof}
Wie schon im Beweis des
Satzes~\ref{buch:holomorph:holomorph:satz:potenzreihe-ableitung}
kann der Konvergenzradius der Reihe $F(z)$ mit der Cauchy-Hadamard-Formel
bestimmt werden.
\index{Cauchy-Hadamard-Formel}%
Er ist
\begin{align*}
\limsup_{n\to\infty} \!\sqrt[n]{\frac1{n+1}|a_n|}
&=
\underbrace{
\lim_{n\to\infty}
\frac{1}{\!\sqrt[n]{n+1}}
}_{\displaystyle =1}
\cdot
\limsup_{n\to\infty}
\!\sqrt[n]{|a_n|}
=
\limsup_{n\to\infty}
\!\sqrt[n]{|a_n|}
=
\frac{1}{\varrho}.
\end{align*}
Die Existenz des Grenzwertes von $1/\!\sqrt[n]{n+1}$ folgt auf die
gleiche Art wie im Beweis von
Satz~\ref{buch:holomorph:holomorph:satz:potenzreihe-ableitung}.
\end{proof}

%
% Die Ableitung der Exponentialfunktion
%
\subsubsection{Die Ableitung der Exponentialfunktion}
Die Exponentialfunktion ist als
\index{Exponentialfunktion}%
\begin{equation}
e^z
=
\exp z
=
\sum_{k=0}^\infty \frac{z^k}{k!}
\label{buch:holomorph:holomorph:eqn:exponentialfunktion}
\end{equation}
definiert worden.
Aus den Ableitungsregeln folgt die formale Ableitung
\[
\frac{d}{dz}
e^z
=
\sum_{k=0}^\infty \frac{d}{dz} \frac{z^k}{k!}
=
\sum_{k=0}^\infty \frac{kz^{k-1}}{k!}
=
\sum_{k=1}^\infty \frac{z^{k-1}}{(k-1)!}
=
\sum_{k=0}^\infty \frac{z^k}{k!}
=
e^z.
\]
Die Ableitung der Exponentialfunktion ist also wieder die
Exponentialfunktion.
Die Exponentialfunktion ist also eine holomorphe Funktion,
die die Differentialgleichung $f'=f$ mit der Anfangsbedingung
$f(0)=1$ erfüllt.

%
% Die trigonometrischen Funktionen
%
\subsubsection{Die trigonometrischen Funktionen}
Für die trigonometrischen Funktionen Sinus und Kosinus hat die
\index{trigonometrische Funktion}%
reelle Analysis die Reihenentwicklungen
\begin{align*}
\cos z
&=
\sum_{k=0}^\infty (-1)^k \frac{z^{2k}}{(2k)!}
=
1 - \frac{z^2}{2!} + \frac{z^4}{4!} - \frac{z^6}{6!} + \frac{z^8}{8!} - \dots
\\
\sin z
&=
\sum_{k=0}^\infty (-1)^k \frac{z^{2k+1}}{(2k+1)!}
=
z - \frac{z^3}{3!} + \frac{z^5}{5!} - \frac{z^7}{7!} + \frac{z^9}{9!} - \dots
\end{align*}
\index{cos@$\cos$}%
\index{sin@$\sin$}%
gefunden.
Sie haben die Ableitungen
\begin{align*}
\frac{d}{dz} \cos z
&=
\sum_{k=0}^{\infty} (-1)^k \frac{d}{dz} \frac{z^{2k}}{(2k)!}
=
\sum_{k=0}^{\infty} (-1)^k \frac{2kz^{2k-1}}{(2k)!}
=
\sum_{k=1}^{\infty} (-1)^k \frac{z^{2k-1}}{(2k-1)!}
\\
&=
-
\sum_{k=1}^{\infty} (-1)^{k-1} \frac{z^{2(k-1)+1}}{(2(k-1)+1)!}
=
-
\sum_{l=1}^{\infty} (-1)^l \frac{z^{2l+1}}{(2l+1)!}
=
-\sin z
\\
\frac{d}{dz} \sin z
&=
\sum_{k=0}^\infty (-1)^k \frac{d}{dz} \frac{z^{2k+1}}{(2k+1)!}
=
\sum_{k=0}^\infty (-1)^k \frac{(2k+1)z^{2k}}{(2k+1)!}
=
\sum_{k=0}^\infty (-1)^k \frac{z^{2k}}{(2k)!}
=
\cos z
\end{align*}
Die aus dem Reellen bekannten Ableitungsregeln der trigonometrischen
Funktionen sind als weiterhing gültig.

Aus den Ableitungsregeln für Sinus und Kosinus folgt mit der Quotientenregel
auch die Ableitungsregeln für den Tangens, seine Ableitung ist
\[
\frac{d}{dz}\tan z
=
\frac{d}{dz}
\frac{\sin z}{\cos z}
=
\frac{\cos^2z +\sin^2 z}{\cos^2z}
=
\frac{1}{\cos^2z}.
\]
\index{tanz@$\tan z$}%
Man beachte, dass der Zähler
\[
Z(z)
=
\cos^2z + \sin^2z
\]
eine holomorphe Funktion ist, deren Ableitungen auf der reellen Achse
verschwinden.
Die Taylorreihe von $Z(z)$ ist also die konstante Funktion $Z(z)=1$.

%
% Die eulersche Formel
%
\subsubsection{Die eulersche Formel}
Die Exponentialfunktion ist durch die Exponentialreihe definiert.

\begin{satz}[Die eulersche Formel]
\index{eulersche Formel}%
\label{buch:holomorph:holomorph:satz:eulerformel}
Für $t\in\mathbb{R}$ gilt
\begin{equation}
e^{it}
=
\cos t + i \sin t.
\label{buch:holomorph:holomorph:eqn:eulerformel}
\end{equation}
\end{satz}

\begin{proof}
Wir setzen $it$ als Argument in die Exponentialreihe ein und erhalten
\begin{align*}
e^{it}
&=
\sum_{k=0}^\infty
\frac{(it)^n}{n!}.
\intertext{Aufgeteilt in gerade und ungerade Potenzen ist diese Summe}
&=
\sum_{k=0}^\infty \frac{(it)^{2k}}{(2k)!}
+
\sum_{k=0}^\infty \frac{(it)^{2k+1}}{(2k+1)!}
\\
&=
\sum_{k=0}^\infty (-1)^k \frac{t^{2k}}{(2k)!}
+
i
\sum_{k=0}^\infty (-1)^k \frac{t^{2k+1}}{(2k+1)!}
=
\cos t + i \sin t.
\end{align*}
Damit ist die eulersche Formel bewiesen.
\end{proof}

\begin{korollar}
$e^{i\pi}+1=0.$
\end{korollar}

Für eine beliebige komplexe Zahl $z=x+iy$ ist also
nach der eulerschen Formel
\begin{equation}
e^z
=
e^{x+iy}
=
e^xe^{iy}
=
e^x(\cos y+i\sin y)
\qquad
\Rightarrow
\qquad
\left\{
\quad
\begin{aligned}
u(x,y) &= e^x \cos y \\
v(x,y) &= e^x \sin y
\end{aligned}
\right.
\label{buch:holomorph:holomorph:eqn:exprealimag}
\end{equation}
Da $e^z$ eine holomorphe Funktion ist, kann die Ableitung allein
durch die Ableitung nach $x$ berechnet werden.
Die Ableitung des Faktors $e^x$ ist aber selbst wieder $e^x$, es folgt
daher
\[
\frac{d}{dz}e^z
=e^x(\cos y+i\sin y)
=
e^z.
\]

Mit der Darstellung
\eqref{buch:holomorph:holomorph:eqn:exprealimag}
von Real- und Imaginärteil der Exponentialfunktion 
können auch sofort die Cauchy-Riemann-Differentialgleichungen
verifiziert werden, es gilt:
\begin{equation}
\left.
\begin{aligned}
\frac{\partial u}{\partial x}
&=
e^x\cos y,
&
\frac{\partial u}{\partial y}
&=
-e^x\sin y,
\\
\frac{\partial v}{\partial x}
&=
e^x \sin y,
&
\frac{\partial v}{\partial y}
&=
e^x \cos y
\end{aligned} 
\quad
\right\}
\quad
\Rightarrow
\quad
\left\{
\quad
\begin{aligned}
\frac{\partial u}{\partial x} &= \phantom{-}\frac{\partial v}{\partial y},
\\
\frac{\partial u}{\partial y} &= -\frac{\partial v}{\partial x}.
\end{aligned}
\right.
\end{equation}
Dass die Cauchy-Riemann-Differentialgleichungen erfüllt sind, bestätigt
erneut, dass die Exponentialfunktion holomorph ist.

%
% Der natürliche Logarithmus einer komplexen Zahl
%
\subsubsection{Der natürliche Logarithmus einer komplexen Zahl}
Wir versuchen, ausgehend von
\eqref{buch:holomorph:holomorph:eqn:exprealimag}
den natürlichen Logarithmus einer komplexen Zahl als Umkehrfunktion
der Exponentialfunktion zu definieren.
Zunächst beachten wir, dass für $z=re^{it}$ folgt, dass
\begin{equation}
e^{z+2\pi i}
=
e^z
e^{2\pi i}
=
e^z(\cos 2\pi + i \sin 2\pi)
=
e^z.
\label{buch:holomorph:holomorph:eqn:expperiodisch}
\end{equation}
Die Exponentialfunktion ist also periodisch mit Periode $2\pi i$.
Der natürliche Logarithmus einer komplexen Funktion kann daher auch nicht
eindeutig bestimmt sein, sondern nur bis auf ein Vielfaches von 
$2\pi i$.

Sei jetzt $z = re^{i\varphi}$ gegeben, wobei $r=|z|$ und
$\varphi=\operatorname{arg}z$.
Gesucht ist eine komplexe Zahl $a+bi$ derart, dass
\[
e^{a+bi}
=
e^a(\cos b + i \sin b)
=
z
=
re^{i\varphi}
\]
gilt.
Daraus kann man ablesen, dass $e^a = r$ oder $a = \log r = \log |z|$
ist und $b = \varphi = \operatorname{arg}z$.

\begin{definition}[natürlicher Logarithmus]
\index{Logarithmus!natürlich}%
\index{naturlicher Logarithmus@natürlicher Logarithmus}%
Der natürlich Logarithmus einer komplexen Zahl ist
\[
\log z
=
\log |z| + i\operatorname{arg}z.
\]
Die komplexen Zahlen der Form
$\log |z| + i\operatorname{arg}z + 2\pi k$ mit $k\in\mathbb{Z}$
sind ebenfalls Logarithmen von $z$.
\index{logz@$\log z$}%
\end{definition}

Aus Satz~\ref{buch:holomorph:holomorph:satz:rechenregeln}
wissen wir bereits, dass der natürliche Logarithmus eine holomorphe
Funktion ist.
Die Ableitung kann mit \eqref{buch:holomoroph:holomorph:eqn:inverseholomorph}
berechnet werden und es gilt
\begin{equation}
\frac{d}{dz}\log z
=
\frac{1}{f'(\log z)}
=
\frac1{z},
\label{buch:holomorph:holomorph:eqn:logz}
\end{equation}
wobei $f(z)=e^z$ die Exponentialfunktion mit $f'(z)=e^z$ ist, was auf
den letzten Ausdruck führt.
Die Gleichung zeigt, dass sich die Logarithmusfunktion statt als
Umkehrfunktion der Exponentialfunktion auch als eine Funktion
mit der Ableitung $1/z$ charakterisieren lässt.

\begin{beispiel}
Die Logarithmusfunktion hat Real- und Imaginärteil
\[
u(x,y)
=
\frac12
\log(x^2+y^2)
\qquad\text{und}\qquad
v(x,y)
=
\arctan\frac{y}{x}.
\]
Wir wollen die Cauchy-Riemann-Differentialgleichungen für diese beiden
Funktion verifizieren.
Dazu müssen die Ableitungen von $u$ und $v$ bestimmt werden, sie sind
\begin{align*}
\frac{\partial u}{\partial x}
&=
\frac12
\frac{2x}{x^2+y^2}
=
\frac{x}{x^2+y^2},
&
\frac{\partial u}{\partial y}
&=
\frac12
\frac{2y}{x^2+y^2}
=
\frac{y}{x^2+y^2},
\\
\frac{\partial v}{\partial x}
&=
\frac{1}{1+(y/x)^2}
\frac{-y}{x^2}
=
\frac{-y}{x^2+y^2},
&
\frac{\partial v}{\partial y}
&=
\frac{1}{1+(y/x)^2}\frac{1}{x}
=
\frac{1}{1+(y/x)^2}\frac{x}{x^2}
=
\frac{x}{x^2+y^2}.
\end{align*}
Die Cauchy-Riemann-Differentialgleichungen sind jetzt offensichtlich.
\end{beispiel}

%
% Algebraische Charakterisierung der Logarithmusfunktion
%
\subsubsection{Algebraische Charakterisierung von Exponential- und Logarithmusfunktion}
Die Charakterisierung \eqref{buch:holomorph:holomorph:eqn:logz}
der Logarithmusfunktion lässt sich noch etwas verallgemeinern, indem
man die Kettenregel bemüht.

\begin{definition}[Logarithmus einer Funktion]
\label{buch:holomorph:holomorph:definition:alglogarithmus}
Sei $K$ ein Funktionenkörper von Funktionen der Variablen $z$ und
$u(z)\in K$ eine Funktion.
Eine \emph{Logarithmus} von $u$ ist eine Funktion $\vartheta(z)\in K$
\index{Logarithmus}%
mit der Eigenschaft
\[
\frac{\partial}{dz}\vartheta(z)
=
\frac{u'(z)}{u(z)}.
\]
\end{definition}

Der Vorteil dieser Definition ist, dass sie weder den Begriff der
Zusammensetzung von Funktionen noch die Umkehrfunktion benötigt.
Die Funktion $\vartheta(z)$ wird nicht als Zusammensetzung der
Logarithmusfunktion mit der Funktion $u(z)$ gewonnen, sondern als
Lösung der Differentialgleichung $\vartheta'(z) = u'(z)/u(z)$.

Auch die Exponentialfunktion ist einer solchen rein algebraischen
Charakterisierung zugänglich.
Die Exponentialfunktion hat die Eigenschaft, dass sie mit ihrer
Ableitung übereinstimmt.
Eine Zusammensetzung mit einer anderen Funktion $u(z)$ hat die
Ableitung
\[
\frac{d}{dz} e^{u(z)}
=
e^{u(z)} u'(z).
\]
Daraus lässt sich die folgende Definition gewinnen, welche wie die
Definition~\ref{buch:holomorph:holomorph:definition:alglogarithmus}
rein algebraisch ist und keinen Bezug auf die Zusammensetzung oder
auch nur auf die gewöhnliche Exponentialfunktion nimmt.

\begin{definition}[Exponentialfunktion, algebraische Charakterisierung]
\label{buch:holomorph:holomorph:definition:algexponential}
Sei $K$ ein Funktionenkörper von Funktionen einer Variable $z$ und $u\in K$.
Eine Funktion $\vartheta\in K$ heisst eine \emph{Exponentialfunktion} von $u$,
\index{Exponentialfunktion}%
wenn
\[
\vartheta'(z) = \vartheta(z)\cdot u'(z)
\]
gilt.
\end{definition}

\begin{beispiel}
Da $K$ ein Funktionenkörper ist, ist die Funktion $u(z)=z^2\in K$.
Die Funktion $e^{z^2}=e^{u(z)}$ erfüllt
\[
\frac{d}{dz} e^{z^2}
=
e^{z^2}\cdot 2z
=
e^{u(z)}\cdot u'(z),
\]
also ist $e^{z^2}$ eine Exponentialfunktion von $u(z)=z^2$.
\end{beispiel}

